\begin{abstract}
Document files with sensitive information are used across nearly every industry. In recent years, cyberattacks targeting cloud-based file transfer applications have resulted in millions of sensitive documents being exposed. Although document encryption methods exist, they are often limited in usability, security, or deployability. Password-based encryption, the most widely available option, remains vulnerable to brute-force attacks, unauthorized sharing, and human error.

We address these gaps in three phases. First, we present a structured comparative framework adapted from the usability-deployability-security model to evaluate nine document encryption methods across 15 design properties, identifying the benefits and limitations of current approaches. Second, we design and implement a novel document encryption system leveraging FIDO2 passkeys, the WebAuthn PRF extension, and open-source age encryption to provide end-to-end protection of digital documents without passwords. Third, we evaluate this system through a counterbalanced within-subjects user study (N=21) comparing the passwordless approach to traditional password-based encryption as provided by Adobe Acrobat.

Mixed-methods results show the passwordless prototype achieved comparable usability to Adobe Acrobat on first use (SUS 72.0 vs.\ 73.3, $p=.944$) while being rated significantly more useful on the Van der Laan acceptance scale ($p=.014$). When asked which system they would prefer for sharing multiple documents, 19 of 21 participants chose the passwordless approach. Qualitative analysis revealed the password sharing burden was a major pain point with the traditional method, while trust in biometric authentication was a key driver of preference for the passwordless system.

To our knowledge, this is the first study to apply FIDO2 and WebAuthn PRF to end-to-end document encryption rather than user authentication, and one of the first human-subjects usability evaluations of the WebAuthn PRF extension in any context.
\end{abstract}
